% Options for packages loaded elsewhere
% Options for packages loaded elsewhere
\PassOptionsToPackage{unicode}{hyperref}
\PassOptionsToPackage{hyphens}{url}
\PassOptionsToPackage{dvipsnames,svgnames,x11names}{xcolor}
%
\documentclass[
  letterpaper,
  DIV=11,
  numbers=noendperiod]{scrartcl}
\usepackage{xcolor}
\usepackage{amsmath,amssymb}
\setcounter{secnumdepth}{5}
\usepackage{iftex}
\ifPDFTeX
  \usepackage[T1]{fontenc}
  \usepackage[utf8]{inputenc}
  \usepackage{textcomp} % provide euro and other symbols
\else % if luatex or xetex
  \usepackage{unicode-math} % this also loads fontspec
  \defaultfontfeatures{Scale=MatchLowercase}
  \defaultfontfeatures[\rmfamily]{Ligatures=TeX,Scale=1}
\fi
\usepackage{lmodern}
\ifPDFTeX\else
  % xetex/luatex font selection
\fi
% Use upquote if available, for straight quotes in verbatim environments
\IfFileExists{upquote.sty}{\usepackage{upquote}}{}
\IfFileExists{microtype.sty}{% use microtype if available
  \usepackage[]{microtype}
  \UseMicrotypeSet[protrusion]{basicmath} % disable protrusion for tt fonts
}{}
\makeatletter
\@ifundefined{KOMAClassName}{% if non-KOMA class
  \IfFileExists{parskip.sty}{%
    \usepackage{parskip}
  }{% else
    \setlength{\parindent}{0pt}
    \setlength{\parskip}{6pt plus 2pt minus 1pt}}
}{% if KOMA class
  \KOMAoptions{parskip=half}}
\makeatother
% Make \paragraph and \subparagraph free-standing
\makeatletter
\ifx\paragraph\undefined\else
  \let\oldparagraph\paragraph
  \renewcommand{\paragraph}{
    \@ifstar
      \xxxParagraphStar
      \xxxParagraphNoStar
  }
  \newcommand{\xxxParagraphStar}[1]{\oldparagraph*{#1}\mbox{}}
  \newcommand{\xxxParagraphNoStar}[1]{\oldparagraph{#1}\mbox{}}
\fi
\ifx\subparagraph\undefined\else
  \let\oldsubparagraph\subparagraph
  \renewcommand{\subparagraph}{
    \@ifstar
      \xxxSubParagraphStar
      \xxxSubParagraphNoStar
  }
  \newcommand{\xxxSubParagraphStar}[1]{\oldsubparagraph*{#1}\mbox{}}
  \newcommand{\xxxSubParagraphNoStar}[1]{\oldsubparagraph{#1}\mbox{}}
\fi
\makeatother

\usepackage{color}
\usepackage{fancyvrb}
\newcommand{\VerbBar}{|}
\newcommand{\VERB}{\Verb[commandchars=\\\{\}]}
\DefineVerbatimEnvironment{Highlighting}{Verbatim}{commandchars=\\\{\}}
% Add ',fontsize=\small' for more characters per line
\usepackage{framed}
\definecolor{shadecolor}{RGB}{241,243,245}
\newenvironment{Shaded}{\begin{snugshade}}{\end{snugshade}}
\newcommand{\AlertTok}[1]{\textcolor[rgb]{0.68,0.00,0.00}{#1}}
\newcommand{\AnnotationTok}[1]{\textcolor[rgb]{0.37,0.37,0.37}{#1}}
\newcommand{\AttributeTok}[1]{\textcolor[rgb]{0.40,0.45,0.13}{#1}}
\newcommand{\BaseNTok}[1]{\textcolor[rgb]{0.68,0.00,0.00}{#1}}
\newcommand{\BuiltInTok}[1]{\textcolor[rgb]{0.00,0.23,0.31}{#1}}
\newcommand{\CharTok}[1]{\textcolor[rgb]{0.13,0.47,0.30}{#1}}
\newcommand{\CommentTok}[1]{\textcolor[rgb]{0.37,0.37,0.37}{#1}}
\newcommand{\CommentVarTok}[1]{\textcolor[rgb]{0.37,0.37,0.37}{\textit{#1}}}
\newcommand{\ConstantTok}[1]{\textcolor[rgb]{0.56,0.35,0.01}{#1}}
\newcommand{\ControlFlowTok}[1]{\textcolor[rgb]{0.00,0.23,0.31}{\textbf{#1}}}
\newcommand{\DataTypeTok}[1]{\textcolor[rgb]{0.68,0.00,0.00}{#1}}
\newcommand{\DecValTok}[1]{\textcolor[rgb]{0.68,0.00,0.00}{#1}}
\newcommand{\DocumentationTok}[1]{\textcolor[rgb]{0.37,0.37,0.37}{\textit{#1}}}
\newcommand{\ErrorTok}[1]{\textcolor[rgb]{0.68,0.00,0.00}{#1}}
\newcommand{\ExtensionTok}[1]{\textcolor[rgb]{0.00,0.23,0.31}{#1}}
\newcommand{\FloatTok}[1]{\textcolor[rgb]{0.68,0.00,0.00}{#1}}
\newcommand{\FunctionTok}[1]{\textcolor[rgb]{0.28,0.35,0.67}{#1}}
\newcommand{\ImportTok}[1]{\textcolor[rgb]{0.00,0.46,0.62}{#1}}
\newcommand{\InformationTok}[1]{\textcolor[rgb]{0.37,0.37,0.37}{#1}}
\newcommand{\KeywordTok}[1]{\textcolor[rgb]{0.00,0.23,0.31}{\textbf{#1}}}
\newcommand{\NormalTok}[1]{\textcolor[rgb]{0.00,0.23,0.31}{#1}}
\newcommand{\OperatorTok}[1]{\textcolor[rgb]{0.37,0.37,0.37}{#1}}
\newcommand{\OtherTok}[1]{\textcolor[rgb]{0.00,0.23,0.31}{#1}}
\newcommand{\PreprocessorTok}[1]{\textcolor[rgb]{0.68,0.00,0.00}{#1}}
\newcommand{\RegionMarkerTok}[1]{\textcolor[rgb]{0.00,0.23,0.31}{#1}}
\newcommand{\SpecialCharTok}[1]{\textcolor[rgb]{0.37,0.37,0.37}{#1}}
\newcommand{\SpecialStringTok}[1]{\textcolor[rgb]{0.13,0.47,0.30}{#1}}
\newcommand{\StringTok}[1]{\textcolor[rgb]{0.13,0.47,0.30}{#1}}
\newcommand{\VariableTok}[1]{\textcolor[rgb]{0.07,0.07,0.07}{#1}}
\newcommand{\VerbatimStringTok}[1]{\textcolor[rgb]{0.13,0.47,0.30}{#1}}
\newcommand{\WarningTok}[1]{\textcolor[rgb]{0.37,0.37,0.37}{\textit{#1}}}

\usepackage{longtable,booktabs,array}
\usepackage{calc} % for calculating minipage widths
% Correct order of tables after \paragraph or \subparagraph
\usepackage{etoolbox}
\makeatletter
\patchcmd\longtable{\par}{\if@noskipsec\mbox{}\fi\par}{}{}
\makeatother
% Allow footnotes in longtable head/foot
\IfFileExists{footnotehyper.sty}{\usepackage{footnotehyper}}{\usepackage{footnote}}
\makesavenoteenv{longtable}
\usepackage{graphicx}
\makeatletter
\newsavebox\pandoc@box
\newcommand*\pandocbounded[1]{% scales image to fit in text height/width
  \sbox\pandoc@box{#1}%
  \Gscale@div\@tempa{\textheight}{\dimexpr\ht\pandoc@box+\dp\pandoc@box\relax}%
  \Gscale@div\@tempb{\linewidth}{\wd\pandoc@box}%
  \ifdim\@tempb\p@<\@tempa\p@\let\@tempa\@tempb\fi% select the smaller of both
  \ifdim\@tempa\p@<\p@\scalebox{\@tempa}{\usebox\pandoc@box}%
  \else\usebox{\pandoc@box}%
  \fi%
}
% Set default figure placement to htbp
\def\fps@figure{htbp}
\makeatother


% definitions for citeproc citations
\NewDocumentCommand\citeproctext{}{}
\NewDocumentCommand\citeproc{mm}{%
  \begingroup\def\citeproctext{#2}\cite{#1}\endgroup}
\makeatletter
 % allow citations to break across lines
 \let\@cite@ofmt\@firstofone
 % avoid brackets around text for \cite:
 \def\@biblabel#1{}
 \def\@cite#1#2{{#1\if@tempswa , #2\fi}}
\makeatother
\newlength{\cslhangindent}
\setlength{\cslhangindent}{1.5em}
\newlength{\csllabelwidth}
\setlength{\csllabelwidth}{3em}
\newenvironment{CSLReferences}[2] % #1 hanging-indent, #2 entry-spacing
 {\begin{list}{}{%
  \setlength{\itemindent}{0pt}
  \setlength{\leftmargin}{0pt}
  \setlength{\parsep}{0pt}
  % turn on hanging indent if param 1 is 1
  \ifodd #1
   \setlength{\leftmargin}{\cslhangindent}
   \setlength{\itemindent}{-1\cslhangindent}
  \fi
  % set entry spacing
  \setlength{\itemsep}{#2\baselineskip}}}
 {\end{list}}
\usepackage{calc}
\newcommand{\CSLBlock}[1]{\hfill\break\parbox[t]{\linewidth}{\strut\ignorespaces#1\strut}}
\newcommand{\CSLLeftMargin}[1]{\parbox[t]{\csllabelwidth}{\strut#1\strut}}
\newcommand{\CSLRightInline}[1]{\parbox[t]{\linewidth - \csllabelwidth}{\strut#1\strut}}
\newcommand{\CSLIndent}[1]{\hspace{\cslhangindent}#1}



\setlength{\emergencystretch}{3em} % prevent overfull lines

\providecommand{\tightlist}{%
  \setlength{\itemsep}{0pt}\setlength{\parskip}{0pt}}



 


\KOMAoption{captions}{tableheading}
\makeatletter
\@ifpackageloaded{caption}{}{\usepackage{caption}}
\AtBeginDocument{%
\ifdefined\contentsname
  \renewcommand*\contentsname{Table of contents}
\else
  \newcommand\contentsname{Table of contents}
\fi
\ifdefined\listfigurename
  \renewcommand*\listfigurename{List of Figures}
\else
  \newcommand\listfigurename{List of Figures}
\fi
\ifdefined\listtablename
  \renewcommand*\listtablename{List of Tables}
\else
  \newcommand\listtablename{List of Tables}
\fi
\ifdefined\figurename
  \renewcommand*\figurename{Figure}
\else
  \newcommand\figurename{Figure}
\fi
\ifdefined\tablename
  \renewcommand*\tablename{Table}
\else
  \newcommand\tablename{Table}
\fi
}
\@ifpackageloaded{float}{}{\usepackage{float}}
\floatstyle{ruled}
\@ifundefined{c@chapter}{\newfloat{codelisting}{h}{lop}}{\newfloat{codelisting}{h}{lop}[chapter]}
\floatname{codelisting}{Listing}
\newcommand*\listoflistings{\listof{codelisting}{List of Listings}}
\makeatother
\makeatletter
\makeatother
\makeatletter
\@ifpackageloaded{caption}{}{\usepackage{caption}}
\@ifpackageloaded{subcaption}{}{\usepackage{subcaption}}
\makeatother
\usepackage{bookmark}
\IfFileExists{xurl.sty}{\usepackage{xurl}}{} % add URL line breaks if available
\urlstyle{same}
\hypersetup{
  pdftitle={Crime and What Else?},
  pdfauthor={Emery Kerani},
  colorlinks=true,
  linkcolor={blue},
  filecolor={Maroon},
  citecolor={Blue},
  urlcolor={Blue},
  pdfcreator={LaTeX via pandoc}}


\title{Crime and What Else?}
\usepackage{etoolbox}
\makeatletter
\providecommand{\subtitle}[1]{% add subtitle to \maketitle
  \apptocmd{\@title}{\par {\large #1 \par}}{}{}
}
\makeatother
\subtitle{Investigating the Predictive Power of Social Factors on
Incarceration Rates}
\author{Emery Kerani}
\date{2026-08-05}
\begin{document}
\maketitle

\renewcommand*\contentsname{Table of contents}
{
\hypersetup{linkcolor=}
\setcounter{tocdepth}{3}
\tableofcontents
}

\section{Title, Team, and
Affiliations}\label{title-team-and-affiliations}

\textbf{Project title:} \emph{Crime and What Else? Investigating the
Predictive Power of Social Factors on Incarceration Rates}

\begin{longtable}[]{@{}lll@{}}
\toprule\noalign{}
Name & Role & Willamette email \\
\midrule\noalign{}
\endhead
\bottomrule\noalign{}
\endlastfoot
Emery Kerani & Owner Product Lead & eskerani@willamette.edu \\
\end{longtable}

\textbf{Peer Stakeholder POs:} Mary Rose Krouse, Shanti Brodnick, Rohan
Srinivas Babu

\textbf{Repository:}
\texttt{https://github.com/eskerani/prison-predictor-capstone/tree/main}

\section{Abstract}\label{abstract}

Crime is often taken for granted as the primary driver of incarceration
in the United States, without accounting for underlying factors
(e.g.~poverty) or large-scale policies that have expanded the
prison-industrial complex to an unprecedented level. This project aims
to evaluate the truth of this dominant narrative by utilizing machine
learning to assess how well a variety of factors are able to predict
incarceration rates. Among both single-feature models and a model using
all features to generate predictions, police spending had the lowest
error rate and highest feature importance, indicating that government
spending on police is more closely associated with incarceration rates
than crime.

\section{Introduction: Problem and Research
Questions}\label{introduction-problem-and-research-questions}

\textbf{Problem statement:} The United States prison system currently
holds nearly 2 million people, and for decades has imprisoned its
inhabitants at one of the highest rates in the world (Lurigio 2024). A
variety of legal factors, including mandatory minimum sentences and
``three strikes'' laws, have contributed to the unprecedented rise in
incarceration in the US since the 1970s and 80s (DeFina and Hannon 2009;
Lurigio 2024). The US incarceration rate has since levelled off,
remaining relatively steady since the 2010s; however, this does not
reflect the nation's crime rate, which has fallen since peaking in in
the 1990s (Pettit and Gutierrez 2018). Given that the US's criminal
legal system wields tremendous power---power that is most often turned
against already disadvantaged individuals and their communities---it is
important that we recognize protective and risk factors for
incarceration, many of which are social in nature (Lurigio 2024; Henry
2020; Pettit and Gutierrez 2018). This project seeks to interrogate the
dominant narrative that crime is the primary cause of incarceration by
examining three other risk factors.

\textbf{Primary research question:} How do the predictive power of
social factors (such as poverty, race, or government spending on police)
compare to that of violent and property crime rates when using support
vector regression to model the national incarceration rate?

\textbf{Changes since M1 proposal:} Given findings that crime rates were
\emph{not} the best predictor of incarceration, the secondary research
question in M1 has been dropped.

\section{Impact and Stakeholders}\label{impact-and-stakeholders}

The key stakeholders in this work are those who are most often targeted
by criminalization and incarceration, primarily poor and BIPOC folks
(Pettit and Gutierrez 2018). Additionally, communities across the
country are continually harmed by increasing incarceration rates, as
incarceration has been shown to negatively impact both community health
and poverty (DeFina and Hannon 2009; Wildeman and Wang 2017). By
critically examining the circumstances that can contribute to
incarceration, this project seeks to uplift the work of organizations
such as Transformative Justice Community and the Oregon Justice Resource
Center, which advocate for solutions to harm that shift our paradigm
away from imprisonment.

My shared theme is social justice and identity-driven disparity
analysis. As others who are interested in social justice but perhaps not
as familiar with the prison-industrial complex as I am, I hope that
they'll point out anywhere I make assumptions/lack context and help
ensure that my project doesn't steer anyone toward misinformed
conclusions.

\section{Related Work and Background}\label{related-work-and-background}

For decades, data science and statistics have played a massively
important role in prison reform and abolition causes; as just one
example, we have known for well over three decades that the population
of Black prisoners is disproportionate to the Black population in the US
(Tonry 1994). Data concerning overall prison populations have also been
key in tracking the effects of legislation concerning crime and
sentencing, many of which contributed to astronomical growth of the US
prison population (Mauer 2001). Police spending is another facet of
prior research, particularly when it comes to questions of how to
effectively and durably address crime---while increased policing may
reduce crime in the short term, studies suggest that diverting money
from policing into social services more reliably reduces crime in the
long term (Beck 2025).

To this point, most of the work applying machine learning algorithms to
data concerning the prison system has focused on predicting or
classifying individual behaviors or risk factors. This project takes a
more abstracted view of the legal system and its contributing factors as
a whole.

\section{Data and Data Engineering}\label{data-and-data-engineering}

This project combines data from five sources, all of which are at the
state level and contain information for the years 1980-2020. Data about
the incarcerated population of each state (excluding federal prisoners)
was obtained from the Bureau of Justice Statistics' Corrections
Statistical Analysis Tool; crime data was downloaded from the FBI's
Summary Reporting System (SRS); data on police spending was collected by
the Census Bureau and made available for download by the Urban
Institute's State and Local Finance Data tool; poverty data was
collected by the Census Bureau and included the population of each
state; and race data was estimated based on demographic information
included in the IPUMS Current Population Survey (CPS). Given the
historical nature of this project, all data sets were ingested as
one-time .csv or .xlsx downloads.

Once downloaded, the datasets were individually cleaned and wrangled
into a standardized format. The IPUMS data was aggregated up to the
state and year level, at which point the datasets could be joined
together using state and year as the key. Using these combined data, all
features were standardized by population (per 100k) based on the
population estimates originally included in the poverty dataset. The
resulting combined dataset contained 2,050 rows and 10 columns: state,
year, total population, incarceration rate, total crime rate, violent
crime rate, property crime rate, poverty rate, the rate of non-white
residents, and police spending rate.

\textbf{Sources:}

\begin{longtable}[]{@{}
  >{\raggedright\arraybackslash}p{(\linewidth - 4\tabcolsep) * \real{0.3636}}
  >{\raggedright\arraybackslash}p{(\linewidth - 4\tabcolsep) * \real{0.2727}}
  >{\raggedright\arraybackslash}p{(\linewidth - 4\tabcolsep) * \real{0.3636}}@{}}
\toprule\noalign{}
\begin{minipage}[b]{\linewidth}\raggedright
Source
\end{minipage} & \begin{minipage}[b]{\linewidth}\raggedright
Role
\end{minipage} & \begin{minipage}[b]{\linewidth}\raggedright
Access
\end{minipage} \\
\midrule\noalign{}
\endhead
\bottomrule\noalign{}
\endlastfoot
Corrections Statistical Analysis Tool & Incarceration counts & One-time
download \\
FBI Summary Reporting System & Crime counts & One-time download \\
Census Bureau via Urban Institute & Police spending data & One-time
download \\
Census Bureau & Historical poverty data & One-time download \\
IPUMS CPS & Race survey data & One-time download \\
\end{longtable}

\textbf{Engineering summary:} Raw data in \texttt{data/raw/}; processed
panel in \texttt{data/processed/}; cleaning script in
\texttt{data/processed/data\_cleaning.qmd}.

\begin{verbatim}
#| label: tbl-data-qc
#| tbl-cap: "Example QC summary (replace with your panel)"
#| echo: true

# Replace with your project data when ready:
# panel <- readRDS(here::here("data/processed/analysis_panel.rds"))
# nrow(panel)
knitr::kable(data.frame(
  check = c("Rows", "Counties", "Date range"),
  status = c("stub", "stub", "stub")
))
\end{verbatim}

\section{Ethics and Responsible Use}\label{ethics-and-responsible-use}

\begin{itemize}
\item
  \textbf{Privacy:} Four out of the five datasets incorporated in this
  project were already highly aggregated to the state and year level.
  Data from the IPUMS CPS was downloaded at the original survey
  granularity (from which identifying household characteristics had been
  removed) and aggregated up to the state/year level. Thus, because
  these data represent a high-level overview, they pose very little risk
  to individuals.
\item
  \textbf{Fairness / equity:} Where associations are found between
  social factors such as race or poverty and incarceration, the data
  must not be interpreted as these factors being indicators of
  criminality. For example, though it has been demonstrated that Black
  people are disproportionately imprisoned in the US, this is due to
  overpolicing and systemic oppression and \emph{not} because of any
  inherent association between Blackness and criminality. Under no
  circumstances should this project be used to justify further policing
  or incarceration of marginalized communities.
\item
  \textbf{Limitations of setting:} The data used in this project come
  with several specific caveats attached. Police spending data was
  missing for all states in 2001 and 2003 through the Urban Institute;
  since the rise in police spending over the years is generally fairly
  steady, these values were imputed by averaging the spending in the
  years directly before and after years with missing data (which allowed
  the preservation of trends by state). For two years (2013 and 2017) of
  poverty data, there were two estimated counts published for both
  population and poverty levels due to new measures being introduced in
  the survey. The new measures were given to only a proportion of the
  population, but because in theory both versions were estimating the
  same thing the values were combined by averaging. Regarding the
  Summary Reporting System's crime data, a number of states were unable
  to provide complete data for various years. Where this happened, the
  FBI generated estimates by either extrapolating using a portion of the
  year's data, or applying regional trends to a state's most recent year
  of complete data. Additionally, several states have changed their
  reporting practices over time, which means that newer numbers may not
  be directly comparable to those collected previously. Finally, the SRS
  only collected information about the most serious offense that
  occurred in a given incident (e.g.~if a homicide occurred in the
  course of a motor vehicle theft, only the homicide would be recorded),
  meaning that these data give us a clearer picture of the number of
  incidents that occurred rather than the actual total number of crimes.
\item
  \textbf{Non-causal framing:} This project only examines the
  associations between variables and thus no conclusions should be made
  about whether or not any of the predictive factors used \emph{cause}
  incarceration. Further, these data and the features used are by nature
  interconnected; some of the features used to generate predictions may
  have impacted others (such as police spending and crime rates).
  Single-variable models in particular should thus be taken with a grain
  of salt as it is impossible to truly isolate any variable.
\end{itemize}

\section{Methods and Evaluation}\label{methods-and-evaluation}

\textbf{Design:} This project uses support vector regression (SVR),
chosen because of its ability to deal with complex, non-linear data
while remaining more straightforward than methods such as XGBoost.
Default parameters were used for all models to ensure that each variable
was treated the same way. Five single-variable models were trained, each
using data from only one predictive factor to ensure that no other data
affected predictions; in addition, one model was trained using all
predictive factors as inputs in order to evaluate the importance of each
factor when combined.

\textbf{Evaluation metrics:} Three main evaluation metrics are used:
correlation coefficient, r2 score, and root mean squared error (RMSE).
The correlation coefficient is used as a preliminary evaluation of the
linearity of relationships between variables, with the expectation that
none will be particularly linear. r2 score is used to evaluate how much
of the variance in incarceration rate can be explained by each model,
while the RMSE measures the average error across predictions.
Additionally, permutation importance and SHAP values are used to
determine the most influential feature of the final multi-variate model.

\textbf{Train / validation / holdout:} An 80-20 split was used to create
training and testing datasets for use with models. This split was
stratified by year to ensure that models were trained on a portion of
the data from each year; the result was that a randomly selected 40
states in each year were used to train models (a total of 1,640 rows),
while the remaining 10 states were used to generate predictions for a
given year (totalling 410 rows). \emph{Only} the relevant data (e.g.~the
feature and national incarceration rate in a single-variable model) were
included in the training and testing data, so that models did not learn
associations between year or state and incarceration rates.
Additionally, all features were scaled using a robust scaler before
training and testing, to eliminate any bias that could by caused by size
differences in the variables. The models' predictions were also scaled
accordingly, but were unscaled before calculating evaluation metrics to
aid interpretibility.

\begin{verbatim}
#| label: fig-methods-placeholder
#| fig-cap: "Replace with a methods diagram or exploratory plot"
#| echo: false
#| fig-width: 5
#| fig-height: 3

plot(1, 1, type = "n", xlab = "", ylab = "", main = "Methods figure placeholder")
text(1, 1, "Replace after M2/M3")
\end{verbatim}

\section{Results and Visualizations}\label{results-and-visualizations}

As expected, none of the features were linearly correlated with state
incarceration rates; all had a ``weak'' correlation coefficient of less
than 0.5. When using single features (i.e.~only police spending) as
predictors of the national incarceration rate, the models generated
performed poorly: only two features (police spending at 89.01 and
property crime rate at 94.41) had RMSEs of less than 100 people per
100k; additionally, the r2 scores (none above 0.3) indicate that no
individual feature was able to explain the variance in the national
incarceration rate particularly well. However, out of these poor models,
the one trained on police spending rates had the best RMSE and r2
scores.

\begin{figure}[H]

{\centering \pandocbounded{\includegraphics[keepaspectratio]{./writeup_files/figure-html/all_graph.png}}

}

\caption{Line plot displaying model predictions of the national
incarceration rate over time using only police spending (in dotted
green) and all features (in dashed orange) as inputs.}

\end{figure}%

A model trained on all features performed markedly better, with an RMSE
of 53.65 and an r2 score of 0.74. Relative to all other features, police
spending rate was the most important in generating this model's
predictions with a permutation importance of 0.16, which shows that the
model's performance suffered the most when police spending data was
shuffled. Similarly, police spending had the highest SHAP value, which
tells us that this was the most influential variable on the model's
performance.

\begin{figure}[H]

{\centering \pandocbounded{\includegraphics[keepaspectratio]{./writeup_files/figure-html/poster_table.png}}

}

\caption{Table containing all model results, with best results
highlighted in dark blue.}

\end{figure}%

Of the decades included in analysis, the multi-variable model generated
the best predictions for 2000 onwward (with RMSEs of 30.41 and 33.5 for
the 2000s and 2010s, respectively) and performed markedly worse on data
from the 1980s (which had an RMSE of 74.48) and 90s (RMSE of 60.51).
With respect to geography, the model's accuracy varied widely: the East
South Central, West North Central, Pacific, and South Atlantic census
regions all had RMSEs of under 50, while other regions ranged from 50 to
71.53. However, it must be noted that predictions were not generated
based on even numbers of instances of states from each region in the
testing data; due to stratification by year when creating the train-test
split, the state distributions were shuffled and thus regional
appearences in testing data vary widely from 27-93.

\section{Discussion, Limitations, and Future
Work}\label{discussion-limitations-and-future-work}

\textbf{Interpretation:} The results above indicate that across the
board, police spending rate serves as the best (or most important)
predictor for national incarceration rate using SVR models. In all
evaluation metrics, property crime rates were a close second, meaning
that where crime was associated with incarceration rates in this project
it was primarily property crime and not violent crime. Additionally, the
performance difference between single-feature and combined-feature
models confirms that none of these variables truly exist in isolation,
and all must be accounted for in order to build a fuller picture of what
might influence incarceration in the United States.

\textbf{Limitations:} One of the chief limitations of this work lies in
the fact that much of the data used were estimates. The data sources
used to calculate race and poverty rates were both surveys, and while
both aim to provide representative information about the US population,
they should not be taken as the definitive truth. Similarly, the data on
crime counts was based only on \emph{reports or arrests made}, which
very likely results in undercounting; out of all reports, only seven
crimes are accounted for in the FBI's Summary Reporting System
(homicide, rape, robbery, aggravated assault, burglary, larceny, and
motor vehicle theft), narrowing the pool even further. Additionally, as
noted previously, this study was purely observational; no causal claims
can be made about the links between any of the variables explored and
incarceration. Rather, the goal of this project is to note associations
between variables.

\textbf{Future work:} Future work could explore variations in the type
of incarceration (e.g.~jails vs.~prisons, men vs.~women) or concentrate
on smaller geographical areas such as counties to incorporate more data.

\section{Conclusions and
Recommendations}\label{conclusions-and-recommendations}

Police spending served as the most influential predictor of national
incarceration rates. Any efforts to reduce incarceration should keep
police and legal system spending in mind as a factor when engineering
reforms.

\section{Reproducibility and
References}\label{reproducibility-and-references}

\textbf{Reproduce main results:}

\begin{Shaded}
\begin{Highlighting}[]
\BuiltInTok{cd}\NormalTok{ your{-}repo}
\ExtensionTok{quarto}\NormalTok{ render deliverables/M4{-}writeup{-}draft/writeup.qmd}
\end{Highlighting}
\end{Shaded}

\textbf{Dependencies:} \texttt{requirements.txt} or \texttt{renv.lock}
in repo root.

\textbf{Data:} URLs and approval status documented in
\texttt{docs/data-dictionary.md}.

\section{References}\label{references}

\phantomsection\label{refs}
\begin{CSLReferences}{1}{0}
\bibitem[\citeproctext]{ref-beck2025}
Beck, Brenden. 2025. {``Local Government Spending: Policing Versus
Social Services.''} \emph{Annual Review of Criminology} 8: 505--28.
\url{https://doi.org/10.1146/annurev-criminol-111523-122639}.

\bibitem[\citeproctext]{ref-defina2013}
DeFina, Robert, and Lance Hannon. 2009. {``The Impact of Mass
Incarceration on Poverty.''} \emph{Crime \& Delinquency} 59 (4):
562--86. \url{https://doi.org/10.1177/0011128708328864}.

\bibitem[\citeproctext]{ref-henry2020}
Henry, Brandy. 2020. {``Adverse Experiences, Mental Health, and
Substance Use Disorders as Social Determinants of Incarceration.''}
\emph{Journal of Community Psychology} 48 (3): 744--62.
\url{https://doi.org/10.1002/jcop.22289}.

\bibitem[\citeproctext]{ref-lurigio2024}
Lurigio, Arthur J. 2024. {``The Golden Anniversary of Mass Incarceration
in America.''} \emph{The Prison Journal} 104 (3): 271--95.
\url{https://doi.org/10.1177/00328855241240131}.

\bibitem[\citeproctext]{ref-mauer2001}
Mauer, Marc. 2001. {``The Causes and Consequences of Prison Growth in
the United States.''} \emph{Punishment \& Society} 3 (1): 9--20.

\bibitem[\citeproctext]{ref-pettit2018}
Pettit, Becky, and Carmen Gutierrez. 2018. {``Mass Incarceration and
Racial Inequality.''} \emph{The American Journal of Economics and
Sociology} 77 (3--4): 1153--82.
\url{https://doi.org/10.1111/ajes.12241}.

\bibitem[\citeproctext]{ref-tonry1994}
Tonry, Michael. 1994. {``Racial Disproportion in US Prisons.''}
\emph{British Journal of Criminology} 34: 97--115.

\bibitem[\citeproctext]{ref-wildeman2017}
Wildeman, Christopher, and Emily A. Wang. 2017. {``Mass Incarceration,
Public Health, and Widening Inequality in the USA.''} \emph{The Lancet}
389 (10077): 1464--74.
\url{https://doi.org/10.1016/S0140-6736(17)30259-3}.

\end{CSLReferences}




\end{document}
